\subsection{The multiplicity-four ring with a rank-35 central datum}
\label{subsec:frontier-629}

The following reduction determines the remaining modular and fusion data
for $R_{6,4,629}$. It does not construct its associators or a category.
All indices in the certificate files are zero based.

\begin{theorem}[Complete necessary modular reduction for $R_{6,4,629}$]
\label{thm:629-modular-reduction}
Every complex categorification of $R_{6,4,629}$ admits a spherical
structure and has a pseudo-unitary Galois conjugate. For that conjugate,
the center has rank $35$ and twist exponent $15$. Its restriction matrix,
normalized modular matrices and fusion tensor belong to one explicitly
computed class, up to complex conjugation and relabelling of simples
having identical restriction columns and twists.
The previously retained induction profiles of ranks $30$ and $38$ are
impossible. The rank-$35$ datum satisfies nonnegative integral Verlinde
multiplicities, but categorifiability remains unresolved.
\end{theorem}

\begin{proof}
The unrestricted signed-sphericalization enumeration has $32$ associative
quotients, all in the sign-gauge orbit of the original tensor. Changing
the chosen lifts by that gauge makes every opposite-lift multiplicity
zero, so the chosen lifts form a tensor subcategory equivalent to the
original category. Its spherical structure transports to the original
category. Of the two rational character components of the original
ring, the quadratic component has zero simple dimensions and cannot be
the spherical dimension character. The other component is one quartic
Galois orbit containing the Frobenius--Perron character. Conjugating to
that embedding gives pseudo-unitarity. The unrestricted enumeration and
the character calculation are exactly replayed by
\filename{research/round4/m4\_629/general\_reduction.json} and its
optimized replay.

The complete induction-profile reduction reconstructs all seven low
frames. Six fail the spectral condition. The remaining frame has two
compatible character labellings and $784$ integral high-column profiles.
Low-word Verlinde integrality and nonnegative restricted products leave
five profiles, of ranks $24,30,32,35,38$. The odd relative-Galois orbit
criterion and exact twist-moment separating certificates reject all five
profiles for one labelling and the ranks $24,32$ for the other. Thus the
input to the present computation is exhaustive: one labelling and three
profiles, with ranks $30,35,38$.

The six low restriction columns, in the retained order, are
\[
\begin{split}
 &(1,0,0,0,0,0),\quad (1,1,1,1,1,1),\quad
 (1,1,2,2,2,3),\\
 &(1,1,2,2,2,2),\quad (1,1,1,2,2,2),\quad
 (1,1,2,1,2,2).
\end{split}
\]
They are all self-dual, as follows from the real character table and its
low-character correspondence. No self-duality assumption is made about
the remaining central simples at this stage.

For each retained profile we rebuild the complete relative root-orbit
domain and the cyclic-rotation trace equations. If a column has
multiplicity $m$ and a relative orbit has degree $e$, its orbit-count
variable lies between $0$ and $\lfloor m/e\rfloor$. Every rotation
eigenmultiplicity lies between zero and the dimension of its rotation
space. These are finite integer intervals. Before rational row reduction
we put variables with larger upper bounds first. This is only a column
permutation, so it preserves the integer lattice and the solution set;
it makes the subsequent exact search much smaller. The surviving free
variables are enumerated with the inherited interval and integrality
pruning, and every leaf is substituted into the original equations.
The complete search results are
\[
\begin{array}{c|rrrrrr}
\text{central rank}&\text{height}&\text{equations}&\text{variables}&
 \text{linear rank}&\text{nodes}&\text{solutions}\\ \hline
30&6&4323&405&350&203&0\\
35&10&2476&587&523&4157&2\\
38&10&1942&527&470&2408&2
\end{array}
\]
Consequently the rank-$30$ profile is excluded. In each other profile the
two solutions are complex conjugate. Every occupied relative orbit has
degree one, so all twists are explicitly in $\mathbb Q(\zeta_{15})$.
The decoded twists have orders $3,5,15$, with least common multiple $15$.
This decoding is exact: a relative average of roots of unity having
absolute value one can only average identical roots. The decoded lists
are also checked to be conjugate as multisets of restriction--twist
pairs. The low simples are exactly the twist-one simples.

Let $D=\operatorname{FPdim}(\mathcal C)$ and write $d_Z$ for a central
simple dimension. The six normalized low rows are fixed by
\[
 S_{A_i,Z}=\lambda_i(Z)d_Z/D,
\]
where $\lambda_i(Z)$ is the joint-character eigenvalue attached to the
restriction column of $Z$. We apply
Proposition~\ref{prop:modular-word-closure}. Since the low objects are
self-dual and exhaust the twist-one objects, the low Galois signed
permutations are determined by unique matches in the known low block.
The congruence identities therefore supply the images of
$T^aSe_i$ for every unit $a$ modulo $15$. Exact modular-word closure
extends the known span from dimension $29$ to dimension $30$ in each
profile. All recorded linear relations, symmetry and isometry identities
are checked in $\mathbb Q(\zeta_{15})$.

For the rank-$38$ profile, the image of the individual basis vector
$e_6$ becomes known. Its restriction is $(0,0,0,0,1,0)$.
The known rows give
\[
 [A_1A_1Z_6:\mathbf1]
 =\sum_j\frac{S_{1j}^{\,2}S_{6j}}{S_{0j}}
 =\frac13,
\]
contradicting Verlinde integrality. The sum is recomputed independently
using rational polynomial arithmetic modulo
\[
 \Phi_{15}(x)=x^8-x^7+x^5-x^4+x^3-x+1.
\]
Conjugating preserves the rational value $1/3$, so both twist branches
are excluded. Only the rank-$35$ profile remains.

For rank $35$, let $W$ be the resulting known $30$-dimensional modular
subspace. Its orthogonal complement is explicitly computed and verified
to be $T$-stable. Its five distinct $T$-eigenvalues are
\[
 \{\zeta_{15}^{2},\zeta_{15}^{3},\zeta_{15}^{8},
   \zeta_{15}^{10},\zeta_{15}^{12}\}.
\]
The center's anomaly-one representation factors through
$SL_2(\mathbb Z/15\mathbb Z)$ by the congruence theorem
\cite[Theorem~II and Remark~4.5]{DLN}. The Chinese remainder theorem
identifies this finite group with
$SL_2(\mathbb F_3)\times SL_2(\mathbb F_5)$.
The computation constructs all $24$ and $120$ matrices in these groups,
their conjugacy classes and the complete class-algebra multiplication
tensors. A class-algebra element with simple spectrum yields every
primitive central character; character orthogonality gives its degree.
Discrete Fourier inversion on powers of the upper unipotent gives its
exact $T$-spectrum. The computation verifies simple spectrum and all
integer degrees and eigenmultiplicities, so no irreducible is omitted.

Tensoring these complete irreducible lists and comparing spectra leaves
exactly two decompositions of $W^\perp$. Their dimensions are
$1+4$ and $1+2+2$. For the first decomposition $-I$ acts trivially; for
the second its negative eigenspace has dimension four. Charge conjugation
can interchange only simples with the same restriction column and twist.
The retained profile has exactly two such pairs, namely
$\{21,22\}$ and $\{23,24\}$, and all other classes are singletons.
An involution on these classes has negative eigenspace dimension at most
two. Thus the $1+2+2$ representation is impossible.

For completeness, the remaining irreducible block is constructed
directly, rather than inferred from approximate character values.
Put $z_p=\zeta_{15}^{15/p}$ and
$g_p=\sum_{x=0}^{p-1}z_p^{x^2}$. Two exact factors are
\[
 S_3=\frac1{g_3}\begin{pmatrix}1&2\\1&-1\end{pmatrix},
 \qquad T_3=\operatorname{diag}(1,z_3),
\]
and
\[
 (S_5)_{xy}=\frac{z_5^{-2xy}-z_5^{2xy}}{g_5}
 \quad(x,y\in\{1,2\}),\qquad
 T_5=\operatorname{diag}(z_5,z_5^4).
\]
Enumerating every Cayley-graph edge checks that these matrices define
representations of the stated finite groups. Their distinct $T$
eigenvalues and nonzero off-diagonal $S$ entries establish irreducibility.
The four-dimensional tensor product has the required spectrum
$\{2,3,8,12\}$ in exponent notation. The remaining one-dimensional
character has $S=1$ and $T=\zeta_{15}^{10}$.

Each complementary $T$-eigenspace is one-dimensional. Hence any other
realization of the surviving representation with this fixed $T$ differs
by a diagonal intertwiner. Symmetry of $S$ fixes the relative squares
of the four diagonal scalars, leaving precisely eight choices of three
signs after removing a common scalar. The reconstruction checks
$S=S^{\mathsf T}$, $S^2=I$ and $STS=T^{-1}ST^{-1}$ for every choice.
Four choices fail nonnegative integral Verlinde multiplicities. The
other four pass all $\binom{36}{3}=7140$ unordered nonunit triples,
and unit triples follow from the unit row and $S^2=I$.
Exact simultaneous row--column permutations show that the four surviving
matrices differ only by the two pair interchanges above. Complex
conjugation supplies the other twist branch. This proves the stated
exhaustive reduction.
\end{proof}

The complete $35\times35$ matrices and $35^3$ fusion tensor are retained
in \filename{research/round5/m4\_629/rank35\_modular\_completions.json}
and \filename{rank35\_fusion\_data.json}. Every nonunit basis object
generates this entire necessary fusion ring. The supplied restriction
matrix is a based-ring homomorphism, verified on all $3780$ unordered
product coordinates. These properties do not supply an associator.

The exact reduction certificate is
\filename{research/round5/m4\_629/modular\_reduction\_certificate.json}.
The script \filename{replay\_modular\_chain.py} reconstructs all three
bounded searches, the twist lists, modular closure, both Verlinde
calculations, the complete finite-group character calculation and every
modular completion. It runs with numerical optimizers disabled and is
also replayed under Python's \texttt{-O} option. The replay records
source and artifact hashes. Numerical infeasibility reports and earlier
search interruptions play no part in the proof.
